% !PRINTED 549
\weberpage{549}
% running head: \S.\,151. Ideale. 549
%<<CONT
ordnen. Dasselbe Ideal $\mathfrak{a}$ ist dann auch sämmtlichen mit $\varphi$
associirten Functionalen zugeordnet. Diese Zuordnung hat die
Eigenschaften 1., 4., 5.

Sind $\varphi$ und $\varphi_1$ zwei nicht associirte Functionale, so ist
gewiss eines von ihnen, etwa $\varphi$, nicht durch das andere $\varphi_1$
theilbar, und folglich giebt es (nach \S.\,143, 3.) ganze Zahlen, die
durch $\varphi$, aber nicht durch $\varphi_1$ theilbar sind. Folglich sind nicht
associirten Functionalen immer \gesperrt{verschiedene} Ideale $\mathfrak{a}$, $\mathfrak{a}_1$ zugeordnet.

Es ist aber nun auch zu zeigen, dass auf diese Weise \gesperrt{alle}
Ideale des Körpers $\Omega$ erhalten werden können, mit anderen
Worten, dass jedes von $\mathfrak{o}$ verschiedene Ideal $\mathfrak{a}$ aus der Gesammtheit
der durch einen gewissen Functionalfactor theilbaren Zahlen
besteht.

Wir gehen also jetzt von irgend einem Ideal $\mathfrak{a}$ aus und
wählen eine beliebige endliche Menge von Zahlen daraus,
$\alpha_1, \alpha_2, \ldots, \alpha_r$, deren grösster gemeinschaftlicher Theiler $\delta_r$ sein
mag. Dieses Functional $\delta_r$ hat eine endliche Anzahl von Primfactoren.

Giebt es nun eine Zahl $\alpha_{r+1}$ in $\mathfrak{a}$, die nicht durch $\delta_r$ theilbar
 ist, so hat der grösste gemeinschaftliche Theiler $\delta_{r+1}$ von
$\delta_r$ und $\alpha_{r+1}$ weniger Primfactoren als $\delta_r$. Wenn wir mit dieser
Schlussweise fortfahren, so kommen wir zu dem Ergebniss, dass
sich aus $\mathfrak{a}$ eine endliche Zahl von Zahlen $\alpha_1, \alpha_2, \ldots, \alpha_m$ so auswählen
 lässt, dass der grösste gemeinschaftliche Theiler $\delta$ dieser
Zahlen in \gesperrt{allen Zahlen von $\mathfrak{a}$ aufgeht.}

Andererseits gehört jede durch $\delta$ theilbare Zahl in $\mathfrak{o}$ zu $\mathfrak{a}$.
Denn nach \S.\,149, 4. kann jede durch $\delta$ theilbare Zahl $\alpha$ in die
Form gesetzt werden

\[
\alpha = \alpha_1 \xi_1 + \alpha_2 \xi_2 + \cdots + \alpha_m \xi_m,
\]

worin $\xi_1, \xi_2, \ldots, \xi_m$ Zahlen in $\mathfrak{o}$ sind, und folglich gehört nach
I. und II. $\alpha$ zum Ideal $\mathfrak{a}$.

Es ergiebt sich hieraus, dass, wenn $\delta$ eine Einheit ist, das
Ideal $\mathfrak{a}$ mit $\mathfrak{o}$ identisch ist. Jedes Ideal $\mathfrak{a}$ ist also dadurch
charakterisirt, dass alle seine Zahlen einen gewissen grössten
gemeinschaftlichen Theiler haben.

Es bleibt noch zu zeigen, dass, wenn die Functionale $\varphi, \psi$
den beiden Idealen $\mathfrak{a}, \mathfrak{b}$ entsprechen, das Product $\varphi\psi$ dem
Ideal $\mathfrak{a}\mathfrak{b}$ entspricht.
